% Generato da docs/62-confronto-crossover-e-verifica-matematica.md con Pandoc.
% Modificare il Markdown sorgente e rigenerare il file TeX.
% Options for packages loaded elsewhere
\PassOptionsToPackage{unicode}{hyperref}
\PassOptionsToPackage{hyphens}{url}
\PassOptionsToPackage{dvipsnames,svgnames,x11names}{xcolor}
\documentclass[
  italian,
  11pt,
  a4paper,
]{article}
\usepackage{xcolor}
\usepackage[margin=22mm]{geometry}
\usepackage{amsmath,amssymb}
\setcounter{secnumdepth}{-\maxdimen} % remove section numbering
\usepackage{iftex}
\ifPDFTeX
  \usepackage[T1]{fontenc}
  \usepackage[utf8]{inputenc}
  \usepackage{textcomp} % provide euro and other symbols
\else % if luatex or xetex
  \usepackage{unicode-math} % this also loads fontspec
  \defaultfontfeatures{Scale=MatchLowercase}
  \defaultfontfeatures[\rmfamily]{Ligatures=TeX,Scale=1}
\fi
\usepackage{lmodern}
\ifPDFTeX\else
  % xetex/luatex font selection
\fi
% Use upquote if available, for straight quotes in verbatim environments
\IfFileExists{upquote.sty}{\usepackage{upquote}}{}
\IfFileExists{microtype.sty}{% use microtype if available
  \usepackage[]{microtype}
  \UseMicrotypeSet[protrusion]{basicmath} % disable protrusion for tt fonts
}{}
\makeatletter
\@ifundefined{KOMAClassName}{% if non-KOMA class
  \IfFileExists{parskip.sty}{%
    \usepackage{parskip}
  }{% else
    \setlength{\parindent}{0pt}
    \setlength{\parskip}{6pt plus 2pt minus 1pt}}
}{% if KOMA class
  \KOMAoptions{parskip=half}}
\makeatother
\usepackage{longtable,booktabs,array}
\usepackage{caption}
\newcounter{none} % for unnumbered tables
\usepackage{calc} % for calculating minipage widths
% Correct order of tables after \paragraph or \subparagraph
\usepackage{etoolbox}
\makeatletter
\patchcmd\longtable{\par}{\if@noskipsec\mbox{}\fi\par}{}{}
\makeatother
\usepackage{graphicx}
\makeatletter
\newsavebox\pandoc@box
\newcommand*\pandocbounded[1]{% scales image to fit in text height/width
  \sbox\pandoc@box{#1}%
  \Gscale@div\@tempa{\textheight}{\dimexpr\ht\pandoc@box+\dp\pandoc@box\relax}%
  \Gscale@div\@tempb{\linewidth}{\wd\pandoc@box}%
  \ifdim\@tempb\p@<\@tempa\p@\let\@tempa\@tempb\fi% select the smaller of both
  \ifdim\@tempa\p@<\p@\scalebox{\@tempa}{\usebox\pandoc@box}%
  \else\usebox{\pandoc@box}%
  \fi%
}
% Set default figure placement to htbp
\def\fps@figure{htbp}
\makeatother
\ifLuaTeX
\usepackage[shorthands=off]{babel}
\else
\usepackage[shorthands=off]{babel}
\fi
\setlength{\emergencystretch}{3em} % prevent overfull lines
\providecommand{\tightlist}{%
  \setlength{\itemsep}{0pt}\setlength{\parskip}{0pt}}
% Stili A4 per i capitoli tecnici. Il kit Altro si legge dalla macchina.
% Sovrascrivere \AltroFontPath prima di caricare questo file per un altro catalogo.
\providecommand{\AltroFontPath}{D:/design-systems/intrawelt-design/assets/fonts/}
\setmainfont{Altro-Grotesk-Book.ttf}[
  Path=\AltroFontPath,
  BoldFont=Altro-Grotesk-Medium.ttf,
  ItalicFont=Altro-Grotesk-Book.ttf,
  ItalicFeatures={FakeSlant=0.16},
  BoldItalicFont=Altro-Grotesk-Medium.ttf,
  BoldItalicFeatures={FakeSlant=0.16}
]
\newfontfamily\intraweltDisplay{Altro-Serif-Light.ttf}[Path=\AltroFontPath,Scale=1]
\setmonofont{Consolas}[Scale=MatchLowercase]

\definecolor{intraweltViolet}{HTML}{4D11A7}
\definecolor{intraweltInk}{HTML}{0A0A0A}
\definecolor{intraweltRule}{HTML}{D0D4E3}
\AtBeginDocument{\color{intraweltInk}}

\usepackage{titlesec}
\usepackage{listings}
\lstset{
  basicstyle=\small\ttfamily,
  breaklines=true,
  breakatwhitespace=true,
  columns=fullflexible,
  keepspaces=true,
  showstringspaces=false
}
\titleformat{\section}
  {\intraweltDisplay\Large\color{intraweltViolet}}{\thesection}{0.7em}{}
\titleformat{\subsection}
  {\intraweltDisplay\large\color{intraweltViolet}}{\thesubsection}{0.7em}{}
\titleformat{\subsubsection}
  {\intraweltDisplay\normalsize\color{intraweltViolet}}{\thesubsubsection}{0.7em}{}
\titlespacing*{\section}{0pt}{2.1ex plus 0.4ex}{0.75ex}
\titlespacing*{\subsection}{0pt}{1.6ex plus 0.3ex}{0.55ex}
\titlespacing*{\subsubsection}{0pt}{1.2ex plus 0.3ex}{0.4ex}

\makeatletter
\renewcommand{\@maketitle}{%
  \begingroup
    \raggedright
    \color{intraweltViolet}\intraweltDisplay
    \fontsize{29}{31}\selectfont \@title\par
  \endgroup
  \vspace{0.65em}%
  {\color{intraweltRule}\hrule height 0.5pt}%
  \vspace{0.75em}%
}
% Booktabs usa già questo hook per il colore dei filetti.
\newcommand*{\CT@arc@}{\color{intraweltRule}}
\makeatother
\AtBeginDocument{%
  \ifdefined\heavyrulewidth
    \setlength{\heavyrulewidth}{0.55pt}%
    \setlength{\lightrulewidth}{0.45pt}%
    \setlength{\cmidrulewidth}{0.45pt}%
  \fi
  \ifdefined\captionsetup
    \captionsetup[figure]{font=small,labelfont=bf}%
  \fi
}

\setlength{\parskip}{0.34\baselineskip plus 1pt}
\setlength{\parindent}{0pt}
\setlength{\emergencystretch}{2em}
\raggedbottom
\widowpenalty=5000
\clubpenalty=5000

\usepackage{bookmark}
\IfFileExists{xurl.sty}{\usepackage{xurl}}{} % add URL line breaks if available
\urlstyle{same}
% fallback for those not using the hyperref driver hyperxmp:
\makeatletter
\@ifundefined{xmpquote}{\newcommand{\xmpquote}[1]{#1}}{}
\makeatother
\hypersetup{
  pdftitle={PC{,} Raspberry Pi e crossover analogico: confronto e verifiche},
  pdflang={it-IT},
  colorlinks=true,
  linkcolor={intraweltViolet},
  filecolor={intraweltViolet},
  citecolor={intraweltViolet},
  urlcolor={intraweltViolet},
  pdfcreator={LaTeX via pandoc}}

\title{PC, Raspberry Pi e crossover analogico: confronto e verifiche}
\author{}
\date{}

\begin{document}
\maketitle

Studio del 2026-11-20, richiesto dall'utente. Il progetto ha scelto il crossover attivo
digitale, ADR-033, e preparato CamillaDSP sul PC Ubuntu già disponibile, ADR-045/046/049.
Con ADR-051 l'utente rende obbligatorio anche l'ascolto con PC spento: si sviluppano
modalità PC e autonoma su Raspberry Pi. La precedente descrizione dell'autonomia come
opzionale è superata. Il \href{sistema-autonomo-driver-e-collaudo.pdf}{progetto aggiornato
di ingressi, driver, finali e collaudo} completa questo confronto; nessun Raspberry o
interfaccia assente viene dichiarato collaudato.

\section{Domanda, metodo e qualità delle
prove}\label{domanda-metodo-e-qualituxe0-delle-prove}

La domanda è se il DSP su PC favorisca una coppia utilizzabile per home recording e ascolto
e quali costi tecnici comporti rispetto a Raspberry Pi o analogico. Il
\href{../../research-vault/crossover-pc-pi-analogico/scope.md}{protocollo congelato prima
dei calcoli} dichiara modello, confronti e tolleranze. Le
\href{../../research-vault/crossover-pc-pi-analogico/sources.md}{fonti S1-S7} sono
documentazione primaria, Linkwitz Lab, nota applicativa e review scientifica. La ricerca è
narrativa e mirata, non sistematica.

Tre livelli di prova rimangono distinti: risultati dedotti dal modello ideale e ricalcolati;
segnali digitali già acquisiti dal nostro banco; misure ancora mancanti di diffusori,
convertitori, stanza, consumi e carico. Le figure qui sono calcoli, non acquisizioni. In
particolare 2 kHz è un esempio, non la scelta dei monitor.

\section{Tre scelte diverse da non sovrapporre}\label{tre-scelte-diverse-da-non-sovrapporre}

La risposta acustica desiderata è un requisito. Il crossover è una funzione di
trasferimento. Analogico e digitale sono modi di realizzarla; PC e Raspberry Pi sono due
possibili calcolatori per la realizzazione digitale. Il numero di finali dipende
dall'architettura attiva: una coppia stereo a due vie richiede quattro canali di potenza, in
entrambi i casi.

Qui l'alternativa analogica è attiva, a livello di linea e prima dei finali. Il
\href{../../docs/52-crossover-passivo-o-attivo.md}{crossover passivo}, con componenti fra un
finale e i due altoparlanti, è un'altra architettura. La soluzione corrente non prevede quel
crossover completo; eventuali protezioni locali del tweeter vanno ancora progettate e non
sostituiscono la funzione acustica da ottenere.

{\def\LTcaptype{none} % do not increment counter
\begin{longtable}[]{@{}
  >{\raggedright\arraybackslash}p{(\linewidth - 6\tabcolsep) * \real{0.2500}}
  >{\raggedright\arraybackslash}p{(\linewidth - 6\tabcolsep) * \real{0.2500}}
  >{\raggedright\arraybackslash}p{(\linewidth - 6\tabcolsep) * \real{0.2500}}
  >{\raggedright\arraybackslash}p{(\linewidth - 6\tabcolsep) * \real{0.2500}}@{}}
\toprule\noalign{}
\begin{minipage}[b]{\linewidth}\raggedright
Aspetto
\end{minipage} & \begin{minipage}[b]{\linewidth}\raggedright
CamillaDSP sul PC esistente
\end{minipage} & \begin{minipage}[b]{\linewidth}\raggedright
CamillaDSP su Raspberry Pi dedicato
\end{minipage} & \begin{minipage}[b]{\linewidth}\raggedright
Crossover analogico attivo
\end{minipage} \\
\midrule\noalign{}
\endhead
\bottomrule\noalign{}
\endlastfoot
Filtro & Parametri software, EQ e ritardi & Stesso modello di filtro, configurazione da
validare sul nuovo host & Reti attive con componenti, guadagni e tolleranze \\
Interfaccia di uscita & Almeno quattro uscite indipendenti per la coppia & DAC/interfaccia
con quattro canali effettivi, da scegliere & Due uscite di linea possono alimentare il
crossover stereo, poi quattro vie analogiche \\
Autonomia dal PC della DAW & PC e sessione audio necessari nella soluzione attuale &
Possibile, se sorgente e percorso d'ingresso del Pi sono davvero autonomi & Possibile con
sorgente analogica e alimentazione del crossover \\
Profili mix/ascolto & Richiamo software pratico e ripetibile & Analogo vantaggio software &
Possibili con EQ, commutazioni o regolazioni aggiuntive \\
Ritardi e fase & Linee di ritardo e filtri programmabili & Stesse possibilità, con limiti di
CPU/I/O da misurare & Correzioni allpass possibili, con compromesso sulla banda \\
Complessità aggiunta & Routing, aggiornamenti, carico condiviso con DAW & Altro host,
alimentazione, storage, I/O e collegamento alla DAW & PCB, operazionali, alimentazione,
taratura e cablaggi \\
Costo iniziale & Riusa il PC; restano interfaccia, finali e protezioni & Aggiunge componenti
rispetto al PC già disponibile & Può riusare una sorgente stereo; aggiunge l'elettronica
analogica \\
Consumo e rumore & Da misurare nel reale uso DAW/ascolto & Da misurare a parità di funzione;
il Pi acceso in più non implica risparmio & Dipendono da elettronica e sorgente effettiva \\
\end{longtable}
}

La documentazione ufficiale CamillaDSP elenca binari x86 e ARM; quella Raspberry Pi
documenta le modalità audio, ma non qualifica il nostro futuro DAC. Nessun prezzo di
acquisto, consumo o numero di xrun dei tre sistemi è stato misurato in questo studio. Le
cifre commerciali storiche in pagina 54 conservano la loro data e non sono un preventivo
aggiornato.

Con un Pi separato va risolto anche il passaggio DAW \(\to\) Pi. Un'interfaccia USB
collegata come periferica a un host non si condivide semplicemente collegandola
contemporaneamente a un secondo host con un cavo sdoppiato. Servirebbero un percorso audio
compatibile e il relativo collaudo, per esempio una trasmissione o un ingresso dedicato; non
sono funzioni già implementate nel progetto. È un requisito di integrazione, non un
impedimento matematico all'uso del Pi.

\section{Dimostrazione 1: lo stesso LR4 esiste in analogico e
digitale}\label{dimostrazione-1-lo-stesso-lr4-esiste-in-analogico-e-digitale}

Nel modello continuo poniamo \(\omega_c=2\pi f_c\). Due sezioni Butterworth del secondo
ordine con \(Q=1/\sqrt{2}\), in cascata, danno:

\[L(s)=\frac{\omega_c^4}{(s^2+\sqrt{2}\omega_cs+\omega_c^2)^2},\qquad H(s)=\frac{s^4}{(s^2+\sqrt{2}\omega_cs+\omega_c^2)^2}.\]

A \(s=j\omega_c\) entrambi valgono -1/2: ciascun ramo è a -6,0205999 dB e la somma coerente,
con la stessa polarità, ha modulo 1. Sommando e fattorizzando:

\[L(s)+H(s)=\frac{s^2-\sqrt{2}\omega_cs+\omega_c^2}{s^2+\sqrt{2}\omega_cs+\omega_c^2}=A(s).\]

Per \(s=j\omega\) il numeratore è il coniugato del denominatore: modulo unitario, fase
variabile. Quindi la somma ideale è allpass, non un segnale privo di sfasamento o un sistema
a fase lineare. Il ritardo di gruppo del modello all'incrocio è:

\[\tau_g(\omega_c)=\frac{2\sqrt{2}}{\omega_c}=0{,}225079\ \mathrm{ms}\quad\text{per }f_c=2000\ \mathrm{Hz}.\]

Questa proprietà è indipendente dal contenitore elettronico che realizza il filtro. Linkwitz
Lab e la review AES sono i riferimenti di contesto; fattorizzazione e valutazione numerica
sono riprodotte dal codice del progetto. Per un diffusore reale occorre ancora la risposta
degli altoparlanti montati.

\begin{figure}
\centering
\pandocbounded{\includegraphics[width=.92\linewidth,keepaspectratio,alt={LR4 ideale: modulo e fase}]{../../build/studio-crossover-20261120-v2/lr4-modulo-fase.pdf}}
\caption{LR4 ideale: modulo e fase}
\end{figure}

Figura prodotta dal nostro modello: due rami, somma coerente e fase della somma. La linea
piatta nel pannello superiore non implica fase costante nel pannello inferiore. Modello
unitario senza driver, stanza o convertitori.

\section{Dimostrazione 2: biquad digitale e frequenza
normalizzata}\label{dimostrazione-2-biquad-digitale-e-frequenza-normalizzata}

Un biquad normalizzato, sezione IIR\footnote{Infinite Impulse Response: filtro ricorsivo la
  cui risposta all'impulso può proseguire indefinitamente, pur decadendo se stabile.} del
secondo ordine, si scrive:

\[y[n]=b_0x[n]+b_1x[n-1]+b_2x[n-2]-a_1y[n-1]-a_2y[n-2].\]

Con \(\omega_0=2\pi f_c/f_s\) e \(\alpha=\sin(\omega_0)/(2Q)\), i coefficienti LPF/HPF del
cookbook W3C si normalizzano dividendo per \(a_0\). Due sezioni identiche per ramo producono
il LR4. A 2 kHz e 48 kHz il calcolo dà:

{\def\LTcaptype{none} % do not increment counter
\begin{longtable}[]{@{}ll@{}}
\toprule\noalign{}
Coefficienti di una sezione & Valori normalizzati \\
\midrule\noalign{}
\endhead
\bottomrule\noalign{}
\endlastfoot
LP \(b_0, b_1, b_2\) & 0,01440144035; 0,02880288069; 0,01440144035 \\
HP \(b_0, b_1, b_2\) & 0,83089802127; -1,66179604255; 0,83089802127 \\
\(a_0, a_1, a_2\) & 1; -1,63299316186; 0,69059892324 \\
\end{longtable}
}

Il massimo modulo dei poli di ciascuna sezione è 0,8310228175, minore di 1. I rami digitali
hanno lo stesso livello di incrocio e somma a modulo unitario del modello ideale. Il
confronto con l'analogico usa la frequenza deformata dalla trasformazione bilineare:

\[f_a=f_c\frac{\tan(\pi f/f_s)}{\tan(\pi f_c/f_s)}.\]

Confrontare direttamente le curve analogica e digitale alla stessa frequenza fino a Nyquist,
ignorando questa trasformazione, sarebbe una verifica sbagliata. Il massimo errore complesso
sul confronto correttamente mappato è circa \(9{,}06\times10^{-15}\). Il calcolo usa float64
NumPy: non è una prova di identità bit per bit fra motori o piattaforme, né qualifica la
quantizzazione di un diverso processore.

\section{Dimostrazione 3: l'analogico richiede Q e guadagno
corretti}\label{dimostrazione-3-lanalogico-richiede-q-e-guadagno-corretti}

Nel Sallen-Key a componenti uguali della nota TI, la frequenza naturale e il Q sono:

\[f_0=\frac{1}{2\pi RC},\qquad Q=\frac{1}{3-K}.\]

Con \(K=1\) si ottiene \(Q=0{,}5\), non Butterworth. Per \(Q=1/\sqrt{2}\) occorre
\(K=3-\sqrt{2}=1{,}5857864\). Due sezioni hanno guadagno \(K^2=2{,}5147186\), circa +8,00979
dB prima della normalizzazione; il progetto analogico deve quindi includere il corretto
bilancio dei livelli. È un esempio della topologia descritta, non una distinta per saldare
il nostro crossover.

Con \(f_c=2\,\mathrm{kHz}\) e \(C=10\,\mathrm{nF}\) risulta \(R=7957{,}747\,\Omega\). Nel
solo modello di scala RC, con R ±1\% e C ±5\%, gli estremi sono:

\[f_{\min}=\frac{2000}{1{,}01\cdot1{,}05}=1885{,}903\ \mathrm{Hz},\qquad f_{\max}=\frac{2000}{0{,}99\cdot0{,}95}=2126{,}528\ \mathrm{Hz}.\]

Sono estremi deterministici di tolleranza, non un intervallo di confidenza. Non descrivono
tutte le variazioni di Q o le correlazioni di quattro componenti reali. Componenti
selezionati, topologie differenti e taratura possono ridurre gli scostamenti. Il digitale
evita la tolleranza RC del calcolo, ma conserva errori di modello, quantizzazione e
interfacce analogiche: nessuna superiorità sonora universale deriva dalla sola etichetta
digitale.

\section{Dimostrazione 4: il filtro elettrico non garantisce la somma
acustica}\label{dimostrazione-4-il-filtro-elettrico-non-garantisce-la-somma-acustica}

Per un lato stereo, in regime lineare, raccogliamo driver montato, guadagni, direttività,
cammino e ambiente nelle funzioni complesse \(G_w\) e \(G_t\):

\[P(f,\theta)=X(f)\left[L(f)G_w(f,\theta)+H(f)G_t(f,\theta)\right].\]

L'uguaglianza del solo filtro è utile ma non impone \(G_w=G_t\). Per mostrare il problema
isoliamo un ritardo puro, mantenendo uguali i livelli dei due contributi al punto di somma.
Una differenza di cammino \(\Delta d\) introduce:

\[\Delta t=\frac{\Delta d}{c},\qquad\Delta\phi=2\pi f\frac{\Delta d}{c},\qquad\left|P(f_c)/X(f_c)\right|=\left|\cos(\Delta\phi/2)\right|.\]

Con \(c=343\,\mathrm{m/s}\) ipotizzata e \(\Delta d=20\,\mathrm{mm}\) si hanno 58,309
microsecondi, 2,798834 campioni a 48 kHz e 41,9825 gradi a 2 kHz: la somma all'incrocio è
-0,596456 dB. Con \(\Delta d=85{,}75\,\mathrm{mm}\) l'errore è mezzo periodo a 2 kHz e la
somma ideale si annulla. Anche invertire la polarità di un ramo del LR4 ideale produce zero
all'incrocio.

Questi sono controesempi calcolati modificando realmente le funzioni complesse. Il caso a
livelli uguali isola il ritardo: non simula da solo la variazione 1/r, la direttività, le
riflessioni o un driver reale. Il DSP può correggere un disallineamento misurato su un asse;
non rende coincidenti due sorgenti in ogni direzione. Non si deduce la risposta della stanza
dai nostri valori elettrici.

\begin{figure}
\centering
\pandocbounded{\includegraphics[width=.92\linewidth,keepaspectratio,alt={Controesempi con ritardo}]{../../build/studio-crossover-20261120-v2/cammino-e-somma.pdf}}
\caption{Controesempi con ritardo}
\end{figure}

Il fondo a -70 dB della figura è un limite di visualizzazione. Non è un rumore misurato. Il
valore al nullo è numericamente circa \(1{,}61\times10^{-16}\) nel modello, non la
previsione di una cancellazione perfetta nella stanza.

\section{Dimostrazione 5: che cosa significa davvero doppio
uso}\label{dimostrazione-5-che-cosa-significa-davvero-doppio-uso}

Se il profilo \(E_p\) è applicato ugualmente ai due rami di una cassa, prima della divisione
o come fattore comune equivalente:

\[P_p=E_p(LG_w+HG_t)X,\qquad\frac{E_pLG_w}{E_pHG_t}=\frac{LG_w}{HG_t}\quad(E_p\ne0).\]

Il rapporto complesso fra i due contributi è preservato. Si può quindi costruire un profilo
mix e uno ascolto sullo stesso crossover di base, senza cambiare arbitrariamente la
relazione tra le vie. Il nostro esempio comune è uno shelving causale dolce; l'errore
relativo calcolato nel rapporto è circa \(6{,}07\times10^{-16}\). La proprietà assume
sistema LTI\footnote{Linear Time-Invariant: sistema lineare e invariabile nel tempo.
  Consente sovrapposizione e descrizione con una funzione di trasferimento fissa.}, stessi
rami e assenza di clipping o interventi dinamici differenti.

Questo dimostra una possibilità architetturale, non che il diffusore rispetti già target,
estensione o SPL desiderati. Un'EQ analogica comune può svolgere una funzione analoga: il
doppio uso non è esclusivo del DSP. Il vantaggio pratico digitale è poter conservare e
richiamare parametri, ritardi e livelli senza rifare reti e tarature meccaniche. I profili
effettivi non sono ancora implementati.

\section{Headroom, escursione, rumore e latenza}\label{headroom-escursione-rumore-e-latenza}

Un aumento di +6 dB moltiplica l'ampiezza per \(10^{6/20}=1{,}99526\). Una sinusoide di
picco 0,8 diventa 1,59621 senza attenuazione preventiva: supera il fondo scala nominale.
Questo esempio non è un limite generale dei picchi di un filtro; per un sistema LTI stabile
vale anche la maggiorazione temporale con la norma L1 dell'impulso,
\(\lVert y\rVert_\infty \le \lVert h\rVert_1\lVert x\rVert_\infty\). Massimo guadagno in
frequenza e massimo picco temporale non sono la stessa cosa.

In un carico resistivo ideale, raddoppiare la tensione quadruplica la potenza. Nel modello
di sorgente piccola a pari pressione, distanza e superficie, l'escursione richiesta varia
come \(1/f^2\): da 60 a 40 Hz aumenta di 2,25 volte. Cabinet, accordo e driver reali
cambiano il limite operativo; l'EQ non crea gratuitamente escursione, dissipazione o volume
di cassa.

Per confrontare il rumore di due convertitori si devono usare stesso riferimento e
banda/pesatura. Solo in quel caso una dinamica DR espressa come rapporto RMS consente
\(V_n=V_\mathrm{ref}\,10^{-DR/20}\). Il guadagno del finale amplifica il rumore d'ingresso e
si aggiunge il rumore proprio. Nessun valore di fruscio al tweeter si ricava dal nome
dell'interfaccia: ADR-046 richiede la prova reale.

Un intervallo di 1024 campioni a 48 kHz dura 21,333 ms; non è la latenza end-to-end già
misurata. Un FIR\footnote{Finite Impulse Response: filtro con risposta all'impulso di
  lunghezza finita; la fase lineare richiede una struttura appropriata dei coefficienti, non
  segue dal solo nome FIR.} causale a fase lineare con \(N=1025\) coefficienti ha ritardo di
gruppo \((N-1)/(2f_s)=10{,}667\,\mathrm{ms}\), cui si aggiunge il resto del percorso. Un IIR
ha fase e ritardo dipendenti dalla frequenza. Anche l'analogico ha ritardo di gruppo, pur
senza buffer digitali. Per registrazione in cuffia si deve progettare il monitoraggio del
musicista separatamente dal percorso dei diffusori.

\section{Riproducibilità e stato delle prove}\label{riproducibilituxe0-e-stato-delle-prove}

Eseguito dalla radice del repository Windows:

\begin{lstlisting}
cd E:\diy-2way-monitors-home
python tools/verifica-matematica-crossover.py --dest build/studio-crossover-20261120-v2
\end{lstlisting}

Il comando richiede una cartella nuova, NumPy, Matplotlib e i font Altro locali indicati
dall'opzione \mbox{\texttt{-\/-font-dir}} o dal kit già presente. Non installa dipendenze e
non tocca motore DSP o hardware. Se la cartella esiste, scegliere un nuovo percorso: nessun
overwrite automatico. Versioni osservate: Python 3.13.14, NumPy 2.2.2, Matplotlib 3.11.0. Il
\href{../../research-vault/crossover-pc-pi-analogico/risultati-modello.json}{rapporto
completo} conserva input, coefficienti, errori, ambiente e hash del sorgente; CSV e formati
PDF/PNG/SVG delle figure sono nei derivati build.

Su 8194 punti deterministici tra 20 e 20000 Hz, incluso \(f_c\) esatto, errore massimo di
modulo LR4 analogico circa \(5{,}56\times10^{-16}\); somma digitale circa
\(9{,}11\times10^{-15}\); uguaglianza con l'allpass circa \(6{,}76\times10^{-16}\). Non sono
8194 esperimenti indipendenti e non si ricava un intervallo statistico. I controesempi di
polarità e ritardo sono verificati separatamente.

Il banco del 19 novembre resta una prova diversa: CamillaDSP reale su grafo virtuale,
acquisizioni, volume, mute e scambio dei lati. Il modello di questo documento spiega perché
aspettarsi certi risultati e quali conclusioni non seguono da essi. Per completare i monitor
servono dati dei driver in cassa, mappa e guadagni delle quattro uscite, livelli e
protezioni, latenza e stabilità sotto DAW, ascolto e misure nella stanza.

\section{Scelta operativa allo stato corrente}\label{scelta-operativa-allo-stato-corrente}

Si sviluppano entrambe le modalità, PC per la DAW e Raspberry Pi per CD/riproduttori a PC
spento, ADR-051. La proposta riusa un'interfaccia USB collegata a un solo host alla volta,
con commutazione da qualificare; in alternativa si sposta il cavo a impianto silenziato. I
filtri matematici sono riutilizzabili, mentre I/O, clock, avvio e latenza richiedono prove
su ciascun host. L'analogico attivo rimane il confronto alternativo. Due crossover completi
in serie richiederebbero un nuovo progetto della funzione complessiva.

La differenza energetica richiede misure. Per un delta di potenza ipotetico \(\Delta P\),
l'energia annua è \(\Delta P\,t/1000\,\mathrm{kWh}\): 5 W o 15 W per 8 ore al giorno
significano 14,6 o 43,8 kWh/anno. Sono scenari, non consumi di PC o Pi misurati. La spesa
totale include elettronica aggiuntiva, integrazione, esercizio e manutenzione: non è
dimostrato che la piattaforma con meno watt nominali sia la più economica per tutti gli usi.

Rimandi: \href{../../docs/61-camilladsp-manuale-tecnico.md}{manuale tecnico},
\href{../../docs/10-ambiente/camilladsp-uso-quotidiano.md}{guida quotidiana},
\href{../../research-vault/crossover-pc-pi-analogico/brief.md}{brief},
\href{../../research-vault/crossover-pc-pi-analogico/review.md}{review e limiti},
\href{../../research-vault/web/linkwitzlab-active-filters.md}{Linkwitz Lab prioritario},
\href{../../docs/INDICE-RAGIONATO.md}{indice generale}.

\end{document}
